% -----------------------------------------------------------------------
% REWRITTEN: Section renamed for EMNLP; observations sharpened with
% linguistic framing.
% -----------------------------------------------------------------------

\section{Question Semantics Predict Visual Information Requirements}\label{sec:background}\label{sec:emp_insights}

\fablenote{For AAAI, consider retitling to ``Why Pre-Retrieval Fidelity Selection Is Both Necessary and Tractable'' --- the current title is NLP-flavored; the proposed one states the two claims the section actually proves and matches a broad-AI audience. Also: all three observations are derived from VQA-v2 + Pixtral-12B with a one-line ``patterns hold across all six models'' assurance. A+ standard is to show the replication, not assert it --- add a small table or appendix pointer with the key statistic (e.g., log-prob AUC, saturation fidelity per question type) for at least 2--3 models.}

We present three empirical observations that collectively establish (i) the practical importance of the pre-retrieval fidelity selection problem and (ii) the tractability of solving it from question text alone.
All observations are based on VQA-v2 with Pixtral-12B across input fidelities ranging from text-only captions to full-resolution images (Figure~\ref{fig:observations_combined}), with qualitatively identical patterns across all six evaluated models.

\textbf{Observation 1: Fidelity Requirements Are Linguistically Predictable.}
Figure~\ref{fig:obs1_dist} shows that query type strongly predicts minimum required fidelity.
Questions containing \emph{how many}, \emph{what color}, or \emph{read/what does it say} signal subtasks requiring fine-grained evidence (object counting, hue discrimination, OCR) that low-fidelity representations destroy; questions containing \emph{is there}, \emph{does the image show}, or \emph{what type} involve coarse semantic judgments resolvable from captions or thumbnails.
On VQA-v2, yes/no and color questions plateau at JPEG Q10 (cost 53.2), while counting, OCR, and spatial queries require full resolution (cost 100).
No single fidelity dominates, establishing that fidelity selection is both necessary and tractable from question text alone~\cite{kafle2017analysis,agrawal2018dont}.
Under-escalation errors (23--25\% of queries, \S\ref{sec:results}) reflect cases where the correct fidelity depends on visual context rather than linguistic structure---the main reducible component of the Oracle--\textsc{VOILA} gap.

\textbf{Observation 2: Post-Retrieval Confidence Is an Unreliable Escalation Signal.}
FrugalGPT-style cascades~\cite{chen2023frugalgpt} retrieve low-cost context first and escalate when VLM log-probability falls below a threshold.
This fails in the fidelity domain: as Figure~\ref{fig:obs2_confidence} shows, incorrect low-fidelity answers receive log-probabilities comparable to---or exceeding---correct ones (43\% of incorrect caption-fidelity answers on VQA-v2 exceed the median log-probability of correct caption-fidelity answers).
For Qwen2-VL-72B, the log-probability AUC for caption-to-full escalation is 0.338---\emph{below random} (ROC curves in Appendix~\ref{sec:ablations})---confirming VLM confidence is an actively misleading escalation signal at low fidelity.
\fablenote{Observation 2 is the paper's strongest novel finding and deserves one sentence of \emph{mechanism}: WHY is confidence anti-correlated at low fidelity? (Plausible: with no visual grounding the VLM falls back on language priors, and prior-driven answers are fluent and hence high-probability.) A causal sentence here turns a curious statistic into an insight reviewers will cite. Also report AUC for all six models, not just the worst one --- cherry-picking the 0.338 invites suspicion that other models are near 0.5.}

\textbf{Observation 3: Model Scaling Does Not Substitute for Sufficient Visual Information.}
Figure~\ref{fig:model_scaling_fidelity} shows that fixing input fidelity at a low level caps accuracy regardless of model size: caption-only accuracy saturates or \emph{degrades} at larger scale (likely due to hallucination when visual grounding is absent), while full-resolution accuracy improves consistently with scale.
Model scaling and information acquisition are therefore \emph{complementary rather than substitutable} optimization axes---the key motivation for pre-retrieval fidelity selection.

These observations establish that fidelity requirements are (1) linguistically predictable pre-retrieval, (2) not estimable from post-hoc confidence, and (3) not compensated by model scaling---motivating \textsc{VOILA}'s design.

\begin{comment}
% OLD OBSERVATION TEXT (kept for reference)
\textbf{Observation 1 (old).}
Figure~\ref{fig:obs1} in the previous section showed that different questions over the same image require different minimum input fidelities. Figure~\ref{fig:obs1_dist} quantifies this at scale: low-fidelity representations (captions, thumbnails) succeed on semantically coarse queries (yes/no, some object identification) but fail disproportionately on counting, color, spatial, and OCR tasks requiring fine-grained visual evidence.
In contrast, high-fidelity inputs substantially improve accuracy on these visually demanding categories while offering limited gains on queries already solvable at low fidelity. No single fidelity dominates across query types, motivating query-adaptive fidelity selection.
\end{comment}
